Der Satellit DFS-3 wurde mittels eines sogenannten \emph{Hohmann-Transfers} auf eine geostationäre Umlaufbahn gebracht: In \hDFO\ Höhe (Perigäum $r_\mathrm{P} = \rpHO$) entliess man den Satelliten auf eine stark exzentrische Übergangsbahn, die im erdfernsten Punkt die geostationäre Bahn schnitt (Apogäum $r_\mathrm{A} = \rAHO$). In diesem Punkt wurden die Antriebsraketen des Satelliten gezündet, um ihm die fehlende Energie für die höherenergetische Kreisbahn zu verleihen.

\begin{abcliste}
\abc Skizzieren Sie die Übergangs- und die geostationäre Bahn des Satelliten und zeichnen Sie die erwähnten Punkte ein.
\abc Die Gesamtenergie eines Satelliten auf einer elliptischen Bahn mit grosser Halbachse $a$ ist (bezüglich eines «unendlich» weit entfernten Punktes)
\begin{align}
 E_\mathrm{G} &= -\frac{1}{2}\,\frac{G\,M_\mathrm{E}\,m_\mathrm{S}}{a}
\end{align}
Berechnen Sie die grosse Halbachse der Übergangsbahn und daraus (mit Hilfe der Energieerhaltung) die Geschwindigkeit des Satelliten im Apogäum vor dem Zünden der Raketen.
\abc Wie gross muss der Geschwindigkeitszuwachs durch die Antriebsraketen sein, damit der Satellit auf die gewünschte geostationäre Kreisbahn gelangt?
\end{abcliste}
